\documentclass[12pt]{scrartcl}

\usepackage{amsmath,amssymb}
\usepackage{fullpage}
\usepackage{hyperref}
\usepackage{graphicx}
\usepackage{tikz}
\usepackage[inline]{enumitem}
\usepackage{tabularx}
\usepackage{array}

\newcolumntype{Y}{>{\centering\arraybackslash}X}
\newcolumntype{L}[1]{>{\raggedright\arraybackslash}p{#1}}
\newcolumntype{C}[1]{>{\centering\arraybackslash}p{#1}}

\setlist[enumerate,1]{label=\alph*), itemjoin=\hspace{1.2em}, itemjoin*=\hspace{1.2em}}

\usetikzlibrary{calc}

\setlength{\parindent}{0pt}

\begin{document}

\begin{center}
	\hrule
	\vspace{0.4cm}
	{\textbf{\large Tutorial 02}}\\ [0.2cm]
    INF1004 Mathematics II
	
\end{center}

\textbf{Name:} Timothy Chia \hspace{\fill} \textbf{Date:} - \\

\hrule

\begin{enumerate}[label=\textbf{\arabic*.}]

\item A spam filter is designed by looking at commonly occurring phrases in spam. Suppose that 80\% of email is spam.
In 10\% of the spam emails, the phrase ``free money'' is used, whereas the phrase ``free money'' is used in only 1\%
of non-spam emails. A new email has just arrived, which does mention ``free money''. What is the probability that it is spam?

\medskip

\textbf{Solution.}
Let $S$ be the event that an email is spam, and let $F$ be the event that an email contains the phrase
\textit{``free money''}.

We are given
\[
P(S)=0.8,\qquad P(F\mid S)=0.1,\qquad P\big(F\mid S^{c}\big)=0.01,
\]
so that $P(S^{c})=1-0.8=0.2$.

We want $P(S\mid F)$, so we first compute $P(F)$ using the Law of Total Probability:
\begin{align*}
P(F)
&=P(F\cap S)+P\big(F\cap S^{c}\big)\\
&=P(F\mid S)P(S)+P\big(F\mid S^{c}\big)P(S^{c})\\
&=0.1(0.8)+0.01(0.2)\\
&=0.08+0.002\\
&=0.082.
\end{align*}
\emph{Note:} We cannot use $P(F\cap S)=P(F)P(S)$ as that would assume independence, which is not given.

Now apply Bayes' Theorem to compute $P(S\mid F)$ in terms of $P(F \mid S)$:
\begin{align*}
P(S\mid F)
&=\frac{P(F\mid S)P(S)}{P(F)}\\
&=\frac{0.1(0.8)}{0.082}\\
&\approx 0.976.
\end{align*}

Therefore, the probability the email is spam given that it contains \textit{``free money''} is approximately
$0.976$ (about $97.6\%$).


%------------------------------------------------------------
\newpage
\item A sample of people, who commute regularly from a town in Surrey into London, were asked for an estimate of the time
taken on their most recent journey. The replies are summarised in the given table.

\medskip

\begin{center}
\setlength{\tabcolsep}{10pt}
\renewcommand{\arraystretch}{1.2}
\begin{tabular}{c|c}
\textbf{Time in minutes} & \textbf{Frequency}\\
\hline
35 to 45  & 12\\
45 to 55  & 54\\
55 to 65  & 68\\
65 to 85  & 41\\
85 to 105 & 23\\
\end{tabular}
\end{center}

\begin{enumerate}[label=(\alph*)]
\item Calculate the estimate of the mean and the standard deviation of these times.
\item What is the median?
\item What is the mode?
\end{enumerate}

\medskip

\textbf{Solution.}

\begin{enumerate}[label=(\alph*)]
\item For grouped data, we estimate using class midpoints. The class midpoints are
\[
40,\ 50,\ 60,\ 75,\ 95
\]
with corresponding frequencies
\[
12,\ 54,\ 68,\ 41,\ 23,
\]
so the total sample size is $n=198$.

\medskip

\textbf{Estimated mean.}
\begin{align*}
\bar{x}
&=\frac{\sum f_i x_i}{\sum f_i}\\
&=\frac{12(40)+54(50)+68(60)+41(75)+23(95)}{198}\\
&=\frac{12520}{198}\\
&\approx 63.23\text{ minutes.}
\end{align*}

\medskip

\textbf{Estimated standard deviation.}
First compute
\[
\sum f_i x_i^2
=12(40^2)+54(50^2)+68(60^2)+41(75^2)+23(95^2)
=837200.
\]
Then
\begin{align*}
s
&\approx \sqrt{\frac{\sum f_i x_i^2}{n}-\bar{x}^{\,2}}\\
&=\sqrt{\frac{837200}{198}-(63.23)^2}\\
&\approx \sqrt{4228.28-3998.33}\\
&\approx \sqrt{229.95}\\
&\approx 15.16\text{ minutes.}
\end{align*}

\item \textbf{Median.}
Compute cumulative frequencies:
\[
12,\ 66,\ 134,\ 175,\ 198.
\]
Since $n=198$, the median position is $n/2=99$, which lies in the class $55$ to $65$.

Let $L=55$ be the lower class boundary, $w=10$ the class width, $c_f=66$ the cumulative frequency before the median class, and $f_m=68$ the frequency of the median class. Then
\begin{align*}
\text{Median}
&= L+\left(\frac{\frac{n}{2}-c_f}{f_m}\right)w\\
&=55+\left(\frac{99-66}{68}\right)10\\
&=55+\frac{33}{68}\cdot 10\\
&\approx 59.85\text{ minutes.}
\end{align*}

\item \textbf{Mode.}
The modal class is $55$ to $65$ (highest frequency $f_1=68$). Let $L=55$, $w=10$, $f_0=54$ (frequency of the class before), and $f_2=41$ (frequency of the class after). Using the grouped-data mode formula,
\begin{align*}
\text{Mode}
&=L+\left(\frac{f_1-f_0}{2f_1-f_0-f_2}\right)w\\
&=55+\left(\frac{68-54}{2(68)-54-41}\right)10\\
&=55+\left(\frac{14}{41}\right)10\\
&\approx 58.41\text{ minutes.}
\end{align*}
\end{enumerate}

%------------------------------------------------------------
\newpage
\item Maurice works at home. At 2 pm he decides to take a break to buy a copy of the \emph{Chronicle} newspaper.
There are three nearby newsagents: Arif, Bob and Carol. However, by 2 pm they may have sold all their \emph{Chronicles} and so have none available.
The independent probabilities that they have a \emph{Chronicle} available at 2 pm are: Arif 0.4, Bob 0.7, Carol 0.25.

\begin{enumerate}[label=(\alph*)]
\item Find the probability that none of the three newsagents have the \emph{Chronicles} at 2 pm.
\item Maurice decides to visit the newsagents in turn until he obtains a \emph{Chronicle} or until he has visited all three.
He tosses a fair coin. If it lands heads he will visit the three newsagents in the order Bob, Arif, Carol.
If it lands tail, he will visit them in order Carol, Arif, Bob.
Find the probability that he will obtain a \emph{Chronicle} from Arif.
\end{enumerate}

\medskip

\textbf{Solution.}
Let $A$, $B$, and $C$ be the events that \emph{Arif}, \emph{Bob}, and \emph{Carol} respectively have a copy of the
\textit{Chronicle} available at 2 pm. We are given
\[
P(A)=0.4,\qquad P(B)=0.7,\qquad P(C)=0.25,
\]
and the events are independent. Hence their complements are
\[
P(A^c)=0.6,\qquad P(B^c)=0.3,\qquad P(C^c)=0.75,
\]
and are also independent.

\begin{enumerate}[label=(\alph*)]
\item The probability that none of the three newsagents have the \textit{Chronicle} at 2 pm is
\begin{align*}
P(A^c\cap B^c\cap C^c)
&=P(A^c)\,P(B^c)\,P(C^c)\\
&=(0.6)(0.3)(0.75)\\
&=0.135.
\end{align*}

\item Let $H$ be the event that the coin lands heads, and $T$ be the event that the coin lands tails.
Then $P(H)=P(T)=0.5$.

\medskip

If $H$ occurs, the visiting order is $B \to A \to C$. Maurice obtains the \textit{Chronicle} from Arif
exactly when Bob does \emph{not} have it and Arif \emph{does} have it, i.e.\ $B^c\cap A$:
\begin{align*}
    P(\text{get from Arif}\mid H) & = P(B^c\cap A) \\
    & = P(B^c)P(A) \\
    & =(0.3)(0.4) \\ 
    & =0.12.
\end{align*}
\medskip
If $T$ occurs, the visiting order is $C \to A \to B$. He obtains the \textit{Chronicle} from Arif
exactly when Carol does \emph{not} have it and Arif \emph{does} have it, i.e.\ $C^c\cap A$:
\begin{align*}
    P(\text{get from Arif}\mid T) & = P(C^c\cap A) \\
    & = P(C^c)P(A) \\
    & =(0.75)(0.4) \\
    & = 0.30.
\end{align*}

Therefore, by the Law of Total Probability,
\begin{align*}
P(\text{get from Arif})
&=P(\text{get from Arif}\mid H)P(H)+P(\text{get from Arif}\mid T)P(T)\\
&=(0.12)(0.5)+(0.30)(0.5)\\
&=0.06+0.15\\
&=0.21.
\end{align*}
\end{enumerate}


%------------------------------------------------------------
\newpage
\item The following table summarises the number of late arrivals of all students who attend a certain school on 5 mornings of particular week.
(Show your steps for all of the following.)

\medskip

\begin{center}
\setlength{\tabcolsep}{14pt}
\renewcommand{\arraystretch}{1.2}
\begin{tabular}{c|c}
\textbf{Number of late arrivals} & \textbf{Number of pupils}\\
\hline
0 & 275\\
1 & 111\\
2 & 33\\
3 & 12\\
4 & 13\\
5 & 16\\
\end{tabular}
\end{center}

\medskip

\begin{enumerate}[label=(\alph*)]
\item Let random variable $X=$ number of late arrivals. Is the random variable discrete or continuous?
\item Draw the probability distribution for $X$. Calculate the mean and the standard deviation of the data in the table.
\item Calculate the expected value.
\end{enumerate}

\medskip

\textbf{Solution.}

\begin{enumerate}[label=(\alph*)]
\item $X$ is \textbf{discrete} because it counts the number of late arrivals (it can only take integer values $0,1,2,\dots$).

\item The total number of pupils is
\[
n = 275+111+33+12+13+16 = 460.
\]
Hence the probability distribution of $X$ is:

\textbf{Solution.}

\begin{enumerate}[label=(\alph*)]
\item $X$ is \textbf{discrete} because it counts the number of late arrivals (it can only take integer values $0,1,2,\dots$).

\item The total number of pupils is
\[
n = 275+111+33+12+13+16 = 460.
\]
Hence the probability distribution of $X$ is:

\begin{center}
\renewcommand{\arraystretch}{2.0}
\setlength{\tabcolsep}{10pt}
\begin{tabular}{c|cccccc}
$x$ & 0 & 1 & 2 & 3 & 4 & 5 \\
\hline
$P(X=x)$
& $\displaystyle \frac{275}{460}$
& $\displaystyle \frac{111}{460}$
& $\displaystyle \frac{33}{460}$
& $\displaystyle \frac{12}{460}$
& $\displaystyle \frac{13}{460}$
& $\displaystyle \frac{16}{460}$ \\
\end{tabular}
\end{center}


\medskip

\textbf{Mean.}
\begin{align*}
\mu = E(X)
&=\sum_x xP(X=x)\\
&=\frac{0(275)+1(111)+2(33)+3(12)+4(13)+5(16)}{460}\\
&=\frac{345}{460}\\
&=0.75.
\end{align*}

\medskip

\textbf{Standard deviation.}
First compute
\begin{align*}
E(X^2)
&=\sum_x x^2P(X=x)\\
&=\frac{0^2(275)+1^2(111)+2^2(33)+3^2(12)+4^2(13)+5^2(16)}{460}\\
&=\frac{959}{460}.
\end{align*}
Then
\begin{align*}
\mathrm{Var}(X)
&=E(X^2)-\big(E(X)\big)^2\\
&=\frac{959}{460}-(0.75)^2\\
&\approx 1.5223,
\end{align*}
so the standard deviation is
\[
\sigma=\sqrt{\mathrm{Var}(X)}\approx \sqrt{1.5223}\approx 1.234.
\]

\item The expected value is
\[
E(X)=0.75,
\]
i.e.\ the expected number of late arrivals is $0.75$.
\end{enumerate}

\medskip

\textbf{Mean.}
\begin{align*}
\mu = E(X)
&=\sum_x xP(X=x)\\
&=\frac{0(275)+1(111)+2(33)+3(12)+4(13)+5(16)}{460}\\
&=\frac{345}{460}\\
&=0.75.
\end{align*}

\medskip

\textbf{Standard deviation.}
First compute
\begin{align*}
E(X^2)
&=\sum_x x^2P(X=x)\\
&=\frac{0^2(275)+1^2(111)+2^2(33)+3^2(12)+4^2(13)+5^2(16)}{460}\\
&=\frac{959}{460}.
\end{align*}
Then
\begin{align*}
\mathrm{Var}(X)
&=E(X^2)-\big(E(X)\big)^2\\
&=\frac{959}{460}-(0.75)^2\\
&\approx 1.5223,
\end{align*}
so the standard deviation is
\[
\sigma=\sqrt{\mathrm{Var}(X)}\approx \sqrt{1.5223}\approx 1.234.
\]

\item The expected value is
\[
E(X)=0.75,
\]
i.e.\ the expected number of late arrivals is $0.75$.
\end{enumerate}


%------------------------------------------------------------
\newpage
\item Terrence walks to school every morning. The probability that he arrives late is 0.15 and independent of whether he arrives late on any other morning.
For a week, in which he decides to walk to school on five mornings, find:

\begin{enumerate}[label=(\alph*)]
\item Probability he arrives late on two or fewer mornings.
\item Probability he arrives late on three or more mornings.
\item Find the mean and standard deviation of the number of mornings on which he arrives late.
\end{enumerate}

\medskip

\textbf{Solution.}
Let $X$ be the number of mornings (out of 5) that Terrence arrives late.
Since the probability of being late on any morning is $p=0.15$ and mornings are independent,
\[
X \sim \mathrm{Bin}(n=5,\,p=0.15).
\]

\begin{enumerate}[label=(\alph*)]
\item
\begin{align*}
P(X\le 2)
&=\sum_{k=0}^{2}\binom{5}{k}(0.15)^k(0.85)^{5-k}\\
&=\binom{5}{0}(0.85)^5+\binom{5}{1}(0.15)(0.85)^4+\binom{5}{2}(0.15)^2(0.85)^3\\
&\approx 0.973388125.
\end{align*}

\item
\begin{align*}
P(X\ge 3)
&=\sum_{k=3}^{5}\binom{5}{k}(0.15)^k(0.85)^{5-k}\\
&=1-P(X\le 2)\\
&\approx 1-0.973388125\\
&\approx 0.026611875.
\end{align*}

\item For a binomial random variable, the mean and variance are
\[
E(X)=np,\qquad \mathrm{Var}(X)=np(1-p).
\]
Hence
\begin{align*}
    E(X) & = 5(0.15) \\ 
         & =0.75
\end{align*}
and
\begin{align*}
    \sigma & = \sqrt{\mathrm{Var}(X)}\\
           & =\sqrt{5(0.15)(0.85)}\\
           & =\sqrt{0.6375}\\
           & \approx 0.798.
\end{align*}
\end{enumerate}

%------------------------------------------------------------
\newpage
\item Bob is a student who travels to and from university by bus. He believes that the probability of having to wait more than six minutes to catch a bus is 0.4
and is independent of the time of day and direction of travel.

\begin{enumerate}[label=(\alph*)]
\item Assuming Bob's beliefs are correct, calculate values for the mean and standard deviation of the number of times he has to wait more than six minutes to catch a bus
in a week when he catches 10 buses.
\item During a thirteen-week period, the number of times (out of 10) he had to wait more than six minutes to catch a bus were as follows:
\[
4,\,8,\,8,\,9,\,3,\,2,\,2,\,7,\,0,\,1,\,5,\,2,\,0
\]
\begin{enumerate}[label=\roman*.]
\item Calculate the mean and standard deviation of this data.
\item Does the answer in (b)i.\ support Bob's beliefs that the probability of having to wait more than six minutes to catch a bus is constant at 0.4 regardless of any factors?
\end{enumerate}
\end{enumerate}

\medskip

\textbf{Solution.}

\begin{enumerate}[label=(\alph*)]
\item Let $X$ be the number of times (out of 10 bus journeys) that Bob has to wait more than 6 minutes.
Assuming Bob's belief is correct and each journey is independent with probability $p=0.4$,
\[
X \sim \mathrm{Bin}(n=10,\,p=0.4).
\]
Hence,
\[
E(X)=np=10(0.4)=4,
\]
and
\[
\mathrm{Var}(X)=np(1-p)=10(0.4)(0.6)=2.4,
\qquad
\text{so } \sigma=\sqrt{2.4}\approx 1.55.
\]

\item
During the 13-week period, the observed values (out of 10) are
\[
4,\,8,\,8,\,9,\,3,\,2,\,2,\,7,\,0,\,1,\,5,\,2,\,0.
\]

\begin{enumerate}[label=\roman*.]
\item The mean is
\[
\bar{x}=\frac{4+8+8+9+3+2+2+7+0+1+5+2+0}{13}
=\frac{51}{13}\approx 3.923.
\]
The (sample) standard deviation is
\[
s=\sqrt{\frac{1}{n-1}\sum_{i=1}^{n}(x_i-\bar{x})^2}.
\]
Here,
\[
\sum_{i=1}^{13}(x_i-\bar{x})^2=\frac{1572}{13},
\]
so
\[
s=\sqrt{\frac{1}{12}\cdot \frac{1572}{13}}
=\sqrt{\frac{131}{13}}
\approx 3.17.
\]

\item Under Bob's model, the expected mean is $4$ and the standard deviation is approximately $1.55$.
The observed mean $\bar{x}\approx 3.923$ is close to $4$, which is broadly consistent with $p\approx 0.4$
(\(\hat{p}=\bar{x}/10\approx 0.392\)). However, the observed standard deviation ($\approx 3.17$) is much larger
than the binomial model's standard deviation ($\approx 1.55$). This suggests the probability may not be constant
from week to week (or the independence assumption may not hold), so the data does \emph{not} strongly support
a constant $p=0.4$ model.
\end{enumerate}
\end{enumerate}

%------------------------------------------------------------
\newpage
\item A shop sells two types of bath cubes: relaxing and invigorating. The owner of a hotel buys 50 cubes from this shop.
For each cube, there is independently a probability of 0.4 that it will be relaxing.

\begin{enumerate}[label=(\alph*)]
\item Find the probability that more than 90\% of the 50 cubes will be relaxing cubes.
\item The 50 cubes included 22 relaxing cubes. Give a reason, whether or not a binomial distribution will provide an appropriate model for the random variable $R$, in each of the following cases:
\begin{enumerate}[label=\roman*.]
\item A guest at the hotel randomly selects one of the 50 cubes. If it is not a relaxing cube, he replaces it and again selects a cube at random.
He continues this procedure until he obtains a relaxing cube. The random variable $R$ denotes the number of cubes he selects until he obtains a relaxing cube.
\item The owner randomly selects 20 of the 50 cubes and places them in a bowl. The random variable $R$ denotes the number of relaxing cubes placed in the bowl; each being evaluated as one trial.
\end{enumerate}
\end{enumerate}

\medskip

\textbf{Solution.}

Let $X$ be the number of relaxing bath cubes among the 50 purchased. Since each cube is relaxing independently with probability $0.4$,
\[
X \sim \mathrm{Bin}(n=50,\,p=0.4).
\]

\begin{enumerate}[label=(\alph*)]
\item “More than $90\%$ of 50” means more than $45$, i.e.\ at least $46$ relaxing cubes. Hence
\begin{align*}
P(X>0.9\times 50)
&=P(X\ge 46)\\
&=\sum_{k=46}^{50}\binom{50}{k}(0.4)^k(0.6)^{50-k}\\
&\approx 1.57\times 10^{-14}.
\end{align*}
So the probability is essentially $0$ (extremely unlikely).

\item The 50 cubes included 22 relaxing cubes.

\begin{enumerate}[label=\roman*.]
\item \textbf{Not binomial.} Here $R$ is the number of selections \emph{until the first relaxing cube is obtained}.
This is not a fixed number of trials, so a binomial model is not appropriate.  
Since the cube is replaced each time a non-relaxing cube is selected, each trial has the same success probability
\[
p=\frac{22}{50}=0.44,
\]
and trials may be treated as independent. Therefore, $R$ is more suitably modelled by a \textbf{geometric distribution} with parameter $p=0.44$.

\item \textbf{Not binomial.} The owner selects 20 cubes from the 50 \emph{without replacement}. The probability of selecting a relaxing cube changes from draw to draw, so the trials are not independent and do not have a constant success probability. A binomial model is therefore not appropriate; the correct model is \textbf{hypergeometric} with $N=50$, $K=22$, and sample size $n=20$.
\end{enumerate}
\end{enumerate}


\end{enumerate}

\end{document}
